%% SYNTHETIC TEST UNIT generated by writing/tools/make_outline_book.py from FINAL_OUTLINE_v2 -- not book text
\chapter{费曼积分的红外结构：钳夹、幂次计数与红外安全}\label{ch:ir}
本单元为测试用的合成文本（由 FINAL\_OUTLINE\_v2 自动生成），陈述内容从略。

\section{第 1 节（测试）}\label{sec:ir-landau}

\begin{definition}[Feynman 参数表示与 Symanzik 多项式]\label{def:ir-landau-setup}
（测试用合成陈述：内容从略。）
\uses{def:prelim-feynman-rules, def:qcd-dimreg}
\end{definition}

\begin{lemma}[包络恒等式]\label{lem:ir-envelope}
（测试用合成陈述：内容从略。）
\uses{def:ir-landau-setup}
\end{lemma}

\begin{theorem}[Landau 条件（参数形式）：解析性与红外有限性判据 \tagR]\label{thm:ir-landau-alpha}
（$R_{\mathrm{PT}}$：在重整化微扰论中对红外正规化的部分子态严格。）
（测试用合成陈述：内容从略。）
\uses{def:ir-landau-setup, lem:ir-envelope}
\end{theorem}

\begin{theorem}[动量空间的钳夹必要条件 \tagR]\label{thm:ir-landau-momentum}
（$R_{\mathrm{PT}}$：在重整化微扰论中对红外正规化的部分子态严格。）
（测试用合成陈述：内容从略。）
\uses{def:ir-landau-setup, lem:appD-gordan}
\end{theorem}

\begin{remark}[Landau 条件的充分性与整体形变 \tagA]\label{rem:ir-landau-status}
（测试用合成陈述：内容从略。）
\uses{thm:ir-landau-momentum, thm:ir-landau-alpha}
\end{remark}
\begin{proofsketch}测试。\end{proofsketch}
\paragraph{地位与缺口}测试：见文献。

\section{第 2 节（测试）}\label{sec:ir-coleman-norton}

\begin{remark}[欧氏区的 Landau 解与红外有限性]\label{rem:ir-euclid-corners}
（测试用合成陈述：内容从略。）
\uses{def:ir-landau-setup, thm:ir-landau-alpha}
\end{remark}

\begin{example}[离壳三角图的 Symanzik 多项式]\label{ex:ir-euclid-triangle}
（测试用合成陈述：内容从略。）
\uses{rem:ir-euclid-corners}
\end{example}

\begin{definition}[钳夹面、内禀坐标与法坐标]\label{def:ir-pinch-surface}
（测试用合成陈述：内容从略。）
\uses{thm:ir-landau-momentum}
\end{definition}

\begin{proposition}[Coleman–Norton 定理（代数形式） \tagR]\label{prop:ir-coleman-norton}
（$R_{\mathrm{PT}}$：在重整化微扰论中对红外正规化的部分子态严格。）
（测试用合成陈述：内容从略。）
\uses{thm:ir-landau-momentum, def:ir-pinch-surface}
\end{proposition}

\begin{proposition}[钳夹面与约化图 \tagA]\label{prop:ir-pinch-classification}
（测试用合成陈述：内容从略。）
\uses{prop:ir-coleman-norton, def:ir-pinch-surface}
\end{proposition}
\begin{proofsketch}测试。\end{proofsketch}
\paragraph{地位与缺口}测试：本陈述的地位由\cref{rem:ir-landau-status}决定。

\section{第 3 节（测试）}\label{sec:ir-pinch}

\begin{definition}[区域标度]\label{def:ir-scalings}
（测试用合成陈述：内容从略。）
\uses{def:prelim-lightfront}
\end{definition}

\begin{proposition}[超表观红外发散度 \tagR]\label{prop:ir-power-counting}
（$R_{\mathrm{PT}}$：在重整化微扰论中对红外正规化的部分子态严格。）
（测试用合成陈述：内容从略。）
\uses{def:ir-scalings, lem:prelim-boost, def:prelim-feynman-rules, prop:appE-degree}
\end{proposition}

\begin{proposition}[超领头构型及其相消 \tagA]\label{prop:ir-super-leading}
（测试用合成陈述：内容从略。）
\uses{prop:ir-power-counting, thm:prelim-slavnov-taylor}
\end{proposition}
\begin{proofsketch}测试。\end{proofsketch}
\paragraph{地位与缺口}测试：本陈述的地位由\cref{asm:appN-envelope}、\cref{asm:appN-ward-product}决定。

\begin{remark}[标度陈述与积分界]\label{rem:ir-scaling-vs-bound}
（测试用合成陈述：内容从略。）
\uses{prop:ir-power-counting}
\end{remark}

\begin{theorem}[领头区域 \tagA]\label{thm:ir-leading-regions}
（测试用合成陈述：内容从略。）
（本命题只涉及部分子态。）
\uses{prop:ir-pinch-classification, prop:ir-power-counting, prop:ir-super-leading, thm:prelim-slavnov-taylor}
\end{theorem}
\begin{proofsketch}测试。\end{proofsketch}
\paragraph{地位与缺口}测试：本陈述的地位由\cref{asm:appN-envelope}、\cref{asm:appN-separation}、\cref{asm:intro-bridge}决定。

\section{第 4 节（测试）}\label{sec:ir-power}

\begin{example}[单圈无质量顶点]\label{ex:ir-vertex}
（测试用合成陈述：内容从略。）
\uses{thm:ir-landau-alpha, prop:ir-coleman-norton, prop:ir-power-counting}
\end{example}

\begin{example}[Feynman 规范下共线区域的标量极化主导]\label{ex:ir-k-dominance}
（测试用合成陈述：内容从略。）
\uses{prop:ir-power-counting}
\end{example}

\begin{example}[$e^+e^-\to q\bar qg$ 的共线区]\label{ex:ir-real-emission}
（测试用合成陈述：内容从略。）
\uses{prop:ir-power-counting, prop:qcd-soft-collinear-limits}
\end{example}

\begin{definition}[Glauber 区]\label{def:ir-glauber}
（测试用合成陈述：内容从略。）
\uses{def:ir-scalings}
\end{definition}

\begin{proposition}[单圈 Glauber 区的钳夹判据 \tagR]\label{prop:ir-glauber-trapping}
（$R_{\mathrm{PT}}$：在重整化微扰论中对红外正规化的部分子态严格。）
（测试用合成陈述：内容从略。）
\uses{def:ir-glauber, thm:ir-landau-momentum}
\end{proposition}

\section{第 5 节（测试）}\label{sec:ir-glauber}

\begin{theorem}[QED 中的 Bloch–Nordsieck 定理 \tagR]\label{thm:ir-bn-qed}
（$R_{\mathrm{PT}}$：在重整化微扰论中对红外正规化的部分子态严格。）
（测试用合成陈述：内容从略。）
% uses: none
\end{theorem}

\begin{theorem}[KLN 定理 \tagA]\label{thm:ir-kln}
（测试用合成陈述：内容从略。）
\uses{prop:prelim-unitarity, thm:ir-bn-qed}
\end{theorem}
\begin{proofsketch}测试。\end{proofsketch}
\paragraph{地位与缺口}测试：见文献。

\begin{proposition}[有色初态时非阿贝尔 Bloch–Nordsieck 抵消的失效 \tagV]\label{prop:ir-bn-failure}
（测试用合成陈述：内容从略。）
\uses{thm:ir-bn-qed, def:prelim-feynman-rules}
\end{proposition}
\paragraph{地位与缺口}测试：本陈述的地位由\cref{rem:dy-dft-premises}决定。

\begin{definition}[红外安全可观测量]\label{def:ir-safe}
（测试用合成陈述：内容从略。）
\uses{def:qcd-ir-safety-informal}
\end{definition}

\begin{theorem}[红外安全可观测量的有限性 \tagA]\label{thm:ir-safe-finite}
（测试用合成陈述：内容从略。）
\uses{def:ir-safe, thm:ir-leading-regions, prop:prelim-unitarity}
\end{theorem}
\begin{proofsketch}测试。\end{proofsketch}
\paragraph{地位与缺口}测试：本陈述的地位由\cref{thm:ir-leading-regions}决定。

\section{第 6 节（测试）}\label{sec:ir-cancellation}

\begin{definition}[区域方法]\label{def:ir-regions}
（测试用合成陈述：内容从略。）
\uses{def:qcd-dimreg, lem:prelim-scaleless}
\end{definition}

\begin{theorem}[欧氏极限的渐近展开 \tagR]\label{thm:ir-regions-euclid}
（$R_{\mathrm{PT}}$：在重整化微扰论中对红外正规化的部分子态严格。）
（测试用合成陈述：内容从略。）
\uses{def:ir-regions}
\end{theorem}

\begin{proposition}[Minkowski 极限下区域方法的地位 \tagP]\label{prop:ir-regions-minkowski}
（测试用合成陈述：内容从略。）
\uses{def:ir-regions, thm:ir-regions-euclid}
\end{proposition}
\paragraph{地位与缺口}测试：见文献。

\begin{remark}[发散区域积分的值依赖于表示]\label{rem:ir-regions-values}
（测试用合成陈述：内容从略。）
\uses{def:ir-regions}
\end{remark}

\begin{example}[一维玩具积分与零区减除]\label{ex:ir-regions-toy}
（测试用合成陈述：内容从略。）
\uses{def:ir-regions, rem:ir-regions-values}
\end{example}
